\subsection{A Brief History of Dynamic Structural Economics}

Dynamic structural economics emerged from the Cowles Commission's groundbreaking work in the 1940s, which established the foundation for relating economic theory to statistical analysis. \citet{Haavelmo1944} introduced the probabilistic framework that distinguished between structural and reduced-form parameters, formalizing the identification problem in simultaneous equations. Koopmans developed order and rank conditions determining which parameters are identifiable given model specifications, while Marschak emphasized the importance of "autonomy" in structural equations---relationships that remain invariant under policy interventions. \citet{HoodKoopmans1953} consolidated these advances into a coherent set of tools, introducing estimation techniques like limited and full-information maximum likelihood. This methodological fusion of economic theory with statistical inference became the defining characteristic of structural economics.

The field advanced significantly through \citet{McFadden1974}'s discrete choice framework, which formalized the connection between utility maximization theory and probabilistic choice behavior. \citet{Heckman1979}'s sample selection model addressed estimation bias from non-random sample selection, treating participation decisions as part of the economic model. These static frameworks laid groundwork for structural estimation in labor economics, industrial organization, and public finance. By the early 1980s, researchers began transitioning from static to dynamic sequential decision problems. \citet{Wolpin1984} modeled fertility decisions over time, and \citet{Miller1984} developed dynamic job matching models.

A milestone was \citet{Rust1987}'s nested fixed-point (NFXP) algorithm, which formulated and estimated a fully specified Markov decision process for engine replacement decisions. For each trial parameter vector, Rust solved the Bellman equation via fixed-point iteration, then maximized the likelihood of observed choices. This methodological breakthrough showed that full-solution estimation of dynamic discrete choice models was computationally feasible, inspiring subsequent innovations to address the "curse of dimensionality." \citet{BerryLevinsohnPakes1995} revolutionized demand estimation with their random coefficients logit model for differentiated products, inverting market share data to recover underlying consumer utilities while allowing flexible substitution patterns with endogenous prices. This "BLP" approach became the workhorse for static demand estimation in industrial organization, later enhanced through techniques integrating micro-data \citep{Petrin2002} and methodological refinements addressing numerical performance \citep{DubeFoxSu2012}.

The computational burden of repeatedly solving dynamic games spurred methodological innovations in the 2000s. \citet{HotzMiller1993} introduced the conditional choice probability (CCP) method, showing that observed choice probabilities could be inverted to recover value function differences without repeatedly solving the dynamic programming problem. \citet{BajariBenkardLevin2007} adapted this insight to dynamic games through a two-step approach---first estimating policy functions from data, then finding structural parameters that rationalize observed behavior. \citet{AguirregabiriaMira2007} proposed nested pseudo-likelihood estimation, iteratively refining two-step estimates to reduce bias. \citet{PesendorferSchmidtDengler2008} developed value function difference methods using equilibrium conditions directly.

Recent advances have further expanded the field's capabilities. \citet{WeintraubBenkardVanRoy2008} introduced oblivious equilibrium concepts to handle high-dimensional state spaces with many agents. \citet{CilibertoTamer2009} developed moment inequality methods addressing equilibrium multiplicity through partial identification. Advances in production function estimation \citep{AckerbergCavesFrazer2015} refined proxy methods for handling endogeneity. The 2010s witnessed integration with machine learning techniques, using neural networks as function approximators for value and policy functions. \citet{FernandezVillaverdeNunoPerla2024} demonstrated how deep learning can mitigate the curse of dimensionality in solving high-dimensional economic models. \citet{MaliarMaliarWinant2021} used neural networks trained via stochastic gradient descent for solving dynamic programming problems, while \citet{AtashbarShi2023} applied deep reinforcement learning to macroeconomic models.
